\documentclass[twocolumn,10pt,letterpaper]{article}
\usepackage[utf8]{inputenc}
\usepackage{amsmath,amssymb,amsfonts}
\usepackage{geometry}
\usepackage{fancyhdr}
\usepackage{braket}

\geometry{margin=0.75in}

\pagestyle{fancy}
\fancyhf{}
\fancyhead[L]{\small PHYSICAL REVIEW A}
\fancyhead[C]{\small VOLUME 51, NUMBER 4}
\fancyhead[R]{\small APRIL 1995}
\fancyfoot[C]{\thepage}

\title{\textbf{\LARGE Quantum coding}}
\author{
  \textbf{Benjamin Schumacher}\\[1ex]
  \textit{Department of Physics, Kenyon College, Gambier, Ohio 43022}
}
\date{\small (Received 26 April 1993; revised manuscript received 29 November 1994)}

\begin{document}
\maketitle

\begin{abstract}
A quantum analogue of Shannon's noiseless coding theorem is formulated and proved. We consider a source producing an ensemble of pure quantum signal states with an associated density operator $\rho$. The von Neumann entropy $S(\rho) = -\mathrm{Tr}(\rho \log_2 \rho)$ represents the minimal number of two-state quantum systems (qubits) per signal required to transmit the quantum states with arbitrarily high fidelity as the block length $N \to \infty$. A subspace typicality theorem is proven for tensor products of Hilbert spaces, showing that the effective dimension of the quantum signal space is $2^{N S(\rho)}$.
\end{abstract}

\section{Introduction}
In classical information theory, Claude Shannon showed in 1948 that the fundamental limit on the noiseless compression of a stationary data source is given by the Shannon entropy:
\begin{equation}
H(X) = -\sum_{i} p_i \log_2 p_i
\end{equation}
Each symbol produced by the source can be encoded into a sequence of binary digits (bits) such that the average length per symbol approaches $H(X)$ as the block length tends to infinity.

The purpose of this paper is to construct the corresponding theory for quantum mechanical communication channels. Rather than transmitting classical sequences of numbers or characters, we consider an information source whose outputs are pure quantum states drawn from a finite or continuous ensemble:
\begin{equation}
\mathcal{E} = \{ (p_i, |\psi_i\rangle) \}
\end{equation}
The ensemble is described macroscopically by the density operator:
\begin{equation}
\rho = \sum_i p_i |\psi_i\rangle \langle \psi_i|
\end{equation}
which is a positive, Hermitian operator on the system's Hilbert space $\mathcal{H}$, with trace $\mathrm{Tr}(\rho) = 1$.

\section{The Qubit as Elementary Information}
In classical systems, the basic unit of information is the binary digit or \textit{bit}. In quantum mechanics, the elementary carrier of quantum information is a two-state quantum system with basis states $\{|0\rangle, |1\rangle\}$. A general state is a linear superposition:
\begin{equation}
|\psi\rangle = \alpha |0\rangle + \beta |1\rangle, \quad |\alpha|^2 + |\beta|^2 = 1
\end{equation}
We refer to such an elementary system as a \textbf{qubit} (quantum bit).

\section{Quantum Ensembles and Typical Subspaces}
Consider a block of $N$ successive signals emitted by the source. The combined state of the block is represented on the tensor product Hilbert space:
\begin{equation}
\mathcal{H}^{\otimes N} = \mathcal{H} \otimes \mathcal{H} \otimes \cdots \otimes \mathcal{H}
\end{equation}
The density operator describing the $N$-block ensemble is the product operator:
\begin{equation}
\rho_N = \rho^{\otimes N} = \rho \otimes \rho \otimes \cdots \otimes \rho
\end{equation}
Let the spectral decomposition of the single-signal density operator be:
\begin{equation}
\rho = \sum_{k=1}^d \lambda_k |e_k\rangle \langle e_k|
\end{equation}
where $\{|e_k\rangle\}$ forms an orthonormal basis of eigenvectors and $\lambda_k$ are the eigenvalues satisfying $\lambda_k \ge 0$ and $\sum_k \lambda_k = 1$. The von Neumann entropy of the source is:
\begin{equation}
S(\rho) = -\mathrm{Tr}(\rho \log_2 \rho) = -\sum_{k=1}^d \lambda_k \log_2 \lambda_k
\end{equation}

For any sequence of indices $\mathbf{k} = (k_1, k_2, \dots, k_N)$, the corresponding product eigenvector in $\mathcal{H}^{\otimes N}$ has probability:
\begin{equation}
P(\mathbf{k}) = \prod_{j=1}^N \lambda_{k_j}
\end{equation}
By the classical law of large numbers, as $N \to \infty$, the quantity $-\frac{1}{N} \log_2 P(\mathbf{k})$ converges in probability to the entropy $S(\rho)$. 

We define the \textbf{$\epsilon$-typical subspace} $\Lambda_\epsilon^{(N)} \subset \mathcal{H}^{\otimes N}$ as the subspace spanned by all eigenvectors $|\mathbf{k}\rangle = |e_{k_1}\rangle \otimes \cdots \otimes |e_{k_N}\rangle$ satisfying:
\begin{equation}
\left| -\frac{1}{N} \log_2 P(\mathbf{k}) - S(\rho) \right| < \epsilon
\end{equation}

\noindent \textbf{Theorem 1 (Subspace Typicality).} For any $\epsilon > 0$ and $\delta > 0$, for sufficiently large $N$:
\begin{enumerate}
\item The total projection weight onto $\Lambda_\epsilon^{(N)}$ satisfies:
\begin{equation}
\mathrm{Tr}(\rho_N \Pi_\epsilon^{(N)}) > 1 - \delta
\end{equation}
where $\Pi_\epsilon^{(N)}$ is the orthogonal projector onto $\Lambda_\epsilon^{(N)}$.
\item The dimension $D_N = \dim \Lambda_\epsilon^{(N)}$ satisfies:
\begin{equation}
(1 - \delta) 2^{N(S(\rho) - \epsilon)} \le D_N \le 2^{N(S(\rho) + \epsilon)}
\end{equation}
\end{enumerate}

\section{The Quantum Noiseless Coding Theorem}
Let Alice wish to transmit an $N$-signal block to Bob through a quantum noiseless channel of $M$ qubits. A block coding scheme of rate $R = M/N$ consists of:
\begin{enumerate}
\item An encoding operation $\mathcal{C}: \mathcal{H}^{\otimes N} \to \mathcal{H}_2^{\otimes M}$, which is a trace-preserving completely positive map or unitary embedding.
\item A decoding operation $\mathcal{D}: \mathcal{H}_2^{\otimes M} \to \mathcal{H}^{\otimes N}$.
\end{enumerate}

The fidelity of the transmission for a pure state $|\Psi\rangle \in \mathcal{H}^{\otimes N}$ is defined by:
\begin{equation}
F(|\Psi\rangle) = \langle \Psi | \mathcal{D}(\mathcal{C}(|\Psi\rangle\langle\Psi|)) | \Psi \rangle
\end{equation}
The ensemble average fidelity is:
\begin{equation}
\bar{F}_N = \sum_{\mathbf{i}} p_{\mathbf{i}} F(|\psi_{i_1}\rangle \otimes \cdots \otimes |\psi_{i_N}\rangle)
\end{equation}

\noindent \textbf{Theorem 2 (Quantum Noiseless Coding Theorem).}
\begin{enumerate}
\item \textit{Direct Part:} If $R > S(\rho)$, then for any $\delta > 0$, there exists an encoding-decoding scheme for sufficiently large $N$ such that:
\begin{equation}
\bar{F}_N > 1 - \delta
\end{equation}
\item \textit{Converse Part:} If $R < S(\rho)$, then as $N \to \infty$, the fidelity $\bar{F}_N \to 0$.
\end{enumerate}

\section{Discussion and Acknowledgments}
The result proves that the von Neumann entropy $S(\rho)$ is not merely an abstract thermodynamic measure, but the operational quantification of quantum information content. Just as Shannon proved that classical bits bound classical communication, $S(\rho)$ qubits bound quantum communication.

The author wishes to thank William K. Wootters for invaluable discussions in which the term ``qubit'' was coined, and Charles H. Bennett for illuminating insights into quantum data compression.

\begin{thebibliography}{99}
\bibitem{shannon} C. E. Shannon, Bell Syst. Tech. J. \textbf{27}, 379 (1948); \textbf{27}, 623 (1948).
\bibitem{neumann} J. von Neumann, \textit{Mathematische Grundlagen der Quantenmechanik} (Springer, Berlin, 1932).
\bibitem{wootters} W. K. Wootters and W. H. Zurek, Nature \textbf{299}, 802 (1982).
\bibitem{bennett} C. H. Bennett and G. Brassard, in \textit{Proc. IEEE Int. Conf. on Computers, Systems and Signal Processing} (1984), p. 175.
\end{thebibliography}

\end{document}
